\chapter{Wie diese Vorlage zu verwenden ist}

\section{Struktur}

\begin{itemize}
    \item \texttt{Diplomarbeit.tex} ist die Hauptdatei, die alle anderen Dateien einbindet.
    \item \texttt{listings-config.tex} enthält die Konfiguration für die Codeblöcke.
    \item \texttt{images/} enthält alle Bilder, die verwendet werden.
    \item \texttt{textparts/} enthält alle Textteile der Diplomarbeit, die eingebunden werden
    \begin{itemize}
        \item \texttt{0/} enthält ein \LaTeX{} Tutorial und diese Erklärung, wie die Vorlage zu verwenden ist. Das sollte irgendwann rausgenommen werden, muss aber nicht gelöscht werden. Es gibt auch eine \texttt{sources.bib} Datei, die nur für das Tutorial ist, und auch rausgenommen werden sollte.
        \item Jede Person die an der Diplomarbeit arbeitet hat ein eigenes Unterverzeichnis, in dem die Kapitel, Quellen und Zeiterfassung gespeichert werden. In der Vorlage sind zwei Unterverzeichnisse, \texttt{1/} und \texttt{2/}, damit man sieht wie das Format mit mehreren Personen aussieht.
        \item \texttt{abstract.tex} ist das Abstrakt oder die Zusammenfassung. Wenn die Diplomarbeit in Englisch geschrieben wird, sollte das Abstrakt auch in Englisch sein.
        \item \texttt{chapter\_assignment.tex} ist die Kapitelzuordnung, die in der Vorlage nur als Beispiel dient. Am Ende wird eingetragen wer welches Kapitel geschrieben hat.
        \item \texttt{conclusion.tex} ist das Fazit. Jeder hat ein eigenes Fazit, für mehr Infos fragt euren Betreuer.
        \item \texttt{danksagungen.tex} ist die Danksagung
        \item \texttt{diplomandenvorstellung.tex} enthält ein Bild und eine kurze Vorstellung von allen Leuten.
        \item \texttt{erklaerung.tex} ist die eidesstaatliche Erklärung, dass die Arbeit selbstständig verfasst wurde. Muss am Ende (auf Papier) unterschrieben werden.
        \item \texttt{objective.tex} enthält eine kurze Beschreibung des Projekts und eine Zielesetzung für jede Person.
        \item \texttt{titel.tex} ist die Titelseite.
    \end{itemize}
\end{itemize}

\section{Empfohlene Vorgehensweise}

\subsection{Am Anfang}

Eine Person sollte als Verantwortlicher für das Dokument festgelegt werden, dass muss nicht der Projektleiter sein, empfohlen sind Personen mit viel Freizeit, Erfahrung mit Git oder einem sehr guten KI-Abo, ich nehme an dass niemand \LaTeX{} kann. Wenn du das bist:
\begin{itemize}
    \item lade dir die Vorlage herunter (wenn nicht schon), solltet ihr das auf Github sehen auch gerne Forken
    \item überlegt euch die Reihenfolge der Personen in der Arbeit, das kann man später noch ändern aber das ist mühsam
    \item schau in textparts/ und erstelle für jede Person ein nummeriertes Unterverzeichnis, 1 und 2 gibt es schon
    \item geh in Diplomarbeit.tex (die Hautptdatei), such "Kapitel einfügen" und gib die neuen Kapitel rein, wenn es die Titel schon gibt auch einfügen
    \item auch in der Hauptdatei bei "Quellen einbinden" die Quellendateien von allen Personen angeben
    \item in \verb|abstract.tex| eine Überschrift für jede Person einfügen
    \item in \verb|conclusion.tex| eine Überschrift für jede Person einfügen
    \item in \verb|danksagungen.tex| das Template für jede Person einfügen (von \verb|\textbf| bis \verb|\newpage|)
    \item in \verb|diplomandenvorstellung.tex| das Template für jede Person kopieren (von \verb|\begin{flushleft}| bis \verb|\end{tabular}| mit \verb|\newpage| dazwischen)
    \item in \verb|erklaerung.tex| die Namen der Personen eintragen
    \item in \verb|objective.tex| eine Überschrift für jede Person einfügen
    \item in \verb|titel.tex| den Titel der Arbeit eintragen, und eure Namen und die der Betreuer
    \item Kommentare tun nicht weh, lasst die Kommentare am besten drin
    \item Die angepasste Vorlage jetzt an alle verteilen. Ich empfehle euch \emph{sehr}, mit Git + VSCode zu arbeiten (siehe \ref{subsec:0-tutorial-vscode}), sollten manche Personen das nicht wollen, gibt es auch eine Anleitung für Overleaf (siehe \ref{subsec:0-tutorial-overleaf}).
\end{itemize}

\subsection{Was jeder zu tun hat}

\begin{itemize}
    \item Jede Person schreibt ihre Kapitel in ihrem Unterverzeichnis, ich empfehle euch, die Dateien zu nummerieren, siehe dieses Beispiel. Eingebunden werden diese Kapitel in der Hauptdatei, unter eurem \verb|\part|
    \item Jeder hat eine eigene Datei für Quellen (um euch Stundenlange merge-Konflikte zu ersparen)
    \item Es gibt nur einen Ordner für Bilder, Unterordner sind aber empfehlenswert.
    \item Jeder schreibt:
    \begin{itemize}
        \item Ein Abstrakt von seiner Diplomarbeit in \verb|abstract.tex|
        \item Ein Fazit in \verb|conclusion.tex|
        \item Eine Danksagung in \verb|danksagungen.tex|
        \item Persönliche Daten und ein Bild in \verb|diplomandenvorstellung.tex|
        \item Eine Zielsetzung (ungefähr halbe Seite) in \verb|objective.tex|
    \end{itemize}
\end{itemize}

Diese Dinge solltet ihr vor der Abgabe auch noch doppelt und dreifach kontrollieren.

\subsection{Am Schluss}

Am Ende ist zu tun:
\begin{itemize}
    \item Die Kapitelzuordnung in \verb|chapter_assignment.tex|.
    \item In \verb|documentation.tex| die Diplomarbeitsdokumentation einfügen, wenn ihr die Arbeit auf Englisch schreibt in beiden Sprachen.
    \item Alles dreimal durchlesen
    \item \emph{Wenn} ihr KI drüberlaufen lasst, \textbf{JEDE. ÄNDERUNG. ÜBERPRÜFEN.}, ist das Sinnvoll, \emph{würde ich das auch so schreiben}, verstehe ich überhautp was da steht. Mit Copilot in VSCode kann man jede Änderung sehen und akzeptieren oder ablehnen. \textbf{LEST JEDE ÄNDERUNG BEVOR IHR SIE AKZEPTIERT}. Wenn die Zeilen in VSCode zu lang werden kann man mit "Alt + Z" den Zeilenumbruch aktivieren.
\end{itemize}

\label{subsec:0-tutorial-vscode}
\subsection{Setup mit VSCode (+ Git)}

Wichtig: eine \LaTeX{} Installation ist notwendig wenn man am lokalen Rechner arbeitet und \emph{mehrere Gigabyte groß}, installieren kann dauern.

Die verantwortliche Person macht ein Git Repository mit der Vorlage, und alle Klonen das. Es ist \emph{nicht notwendig}, auf eigenen Branches arbeiten, die Vorlage sollte es grundsätzlich aushalten, falsch ist das klassische Git Schema (Branch, Pull Request, Merge) aber sicher nicht.

\subsubsection{Windows}
Ladet euch \LaTeX{} über \href{https://www.tug.org/texlive/acquire-netinstall.html}{TeX Live} herunter, wenn ihr gefragt werdet welche Pakete ihr installieren wollt, nehmt "alle" oder "vollständig".

Wenn ihr die Vorlage in VSCode aufmacht, bekommt ihr Extensions empfohlen, die installieren.

\subsubsection{Linux}
Es wird angenommen, dass VSCode installiert ist.

\begin{lstlisting}[caption={Latex mit apt}]
sudo apt install texlive-full
\end{lstlisting}

\begin{lstlisting}[caption={Latex mit pacman}]
sudo pacman -S texlive-full
\end{lstlisting}

\begin{lstlisting}[caption={VSCode Extensions}]
code --install-extension tecosaur.latex-utilities
code --install-extension james-yu.latex-workshop
code --install-extension mathematic.vscode-latex
code --install-extension christian-kohler.path-intellisens
\end{lstlisting}

\subsubsection{In VSCode wichtig}
\begin{itemize}
    \item Im Code "Strg + Alt + J" öffnet das Dokument an der Stell vom Cursor.
    \item andersherum, im Dokument "Strg + Klick" öffnet die Stelle im Code.
    \item Der \TeX-Reiter an der linken Seite hat nützliche Kommands, eine Übersicht über das Dokument, und Mathe-Symbole
\end{itemize}

\label{subsec:0-tutorial-overleaf}
\subsection{Setup mit Overleaf}
Die Verantwortliche Person erstellt ein Zip-Archiv der Vorlage und schickt es den Personen, die auf Overleaf arbeiten. Overleaf erlaubt Zusammenarbeit im Gratis-Tier nur für zwei Personen, ist also dafür nicht gedacht.

In Overleaf gibt es Buttons zum wechseln zwischen Code und Dokument, alternativ auch Doppelklich im Dokument, um die Stelle im Code zu öffnen, und "Strg + Enter" um neuzukompilieren.

Ich empfehle, nur die eigentliche Arbeit in Overleaf zu schreiben und die anderen Teile (Abstrakt, Fazit, etc.) später in VSCode einzufügen. Dazu sollte auch kein installiertes \LaTeX{} notwendig sein.

Wenn die Arbeit von Overleaf zusammengeführ werden soll, muss theoritisch nur das nummerierte Unterverzeichnis eingefügt werden, Kapitel und Quellen eingebunden, und die anderen Teile ergänzt werden.

\section{Sonstiges}

\makeatletter
Sollte es Fragen oder Verbesserungsvorschläge an der Vorlage geben, könnt ihr mich unter \href{mailto:christian.patzl@gmx.net}{christian.patzl@gmx.net} anschreiben, erwartet nicht, dass ich in einer Woche antworte. Andernfalls sucht jemanden, der jemanden kennt, der mich noch kennt, oder erstellt ein Issue auf \href{https://github.com/Christiano300/DiplVorlage}{Github}.
\makeatother